% ATTEMPT_STATUS: unresolved
% ATTEMPT_ORIGIN: new research, 2026-10-01
% TOP_PROBLEM: {"id": 1441, "title": "Second Neighborhood Problem", "category_id": 3, "category": {"id": 3, "name": "graph_theory", "display_name": "Graph Theory", "description": "Problems involving graphs, networks, and their properties.", "slug": "graph-theory", "order_index": 3, "created_at": "2026-07-31T15:26:25.671Z"}, "status": "open"}
% FIRST_POSTED: 2026-10-01
% Imported from: unsolvedmath_additional_500.tex; UnsolvedMath ID 1441
% !TeX program = lualatex
% Compile twice: lualatex 1441.tex
% Self-contained: no external bibliography, images, or \input files.
% Numeric section labels and visible numbers are original UnsolvedMath IDs.
\documentclass[11pt,a4paper]{article}
\usepackage[margin=24mm,headheight=14pt]{geometry}
\usepackage{fontspec}
\setmainfont{Latin Modern Roman}
\setsansfont{Latin Modern Sans}
\setmonofont{Latin Modern Mono}
\usepackage{amsmath,amssymb,mathtools}
\usepackage{unicode-math}
\setmathfont{Latin Modern Math}
\usepackage{microtype}
\usepackage{enumitem,xurl,needspace}
\usepackage[dvipsnames]{xcolor}
\usepackage{hyperref}
\hypersetup{unicode=true,colorlinks=true,linkcolor=MidnightBlue,urlcolor=MidnightBlue,
 pdftitle={UnsolvedMath Problem 1441},pdfauthor={Source-annotated compilation},bookmarksnumbered=true}
\usepackage{bookmark,tocloft,titlesec,fancyhdr}
\setlength{\cftsecnumwidth}{4.2em}
\cftsetpnumwidth{2.5em}
\cftsetrmarg{4em}
\titleformat{\section}{\Large\bfseries\raggedright}{\thesection}{0.7em}{}
\pagestyle{fancy}\fancyhf{}
\fancyhead[L]{\small\sffamily UnsolvedMath research attempt}
\fancyhead[R]{\small Original catalogue IDs}
\fancyfoot[C]{\thepage}
\setlength{\parindent}{0pt}
\setlength{\parskip}{5pt plus 1pt}
\setlength{\emergencystretch}{3em}
\setcounter{tocdepth}{1}\setcounter{secnumdepth}{1}
\setlist[itemize]{leftmargin=3em,itemsep=4pt,topsep=4pt,parsep=0pt}
\allowdisplaybreaks
\newcommand{\N}{\mathbb{N}}
\newcommand{\Z}{\mathbb{Z}}
\newcommand{\Q}{\mathbb{Q}}
\newcommand{\R}{\mathbb{R}}
\newcommand{\C}{\mathbb{C}}
\newcommand{\F}{\mathbb{F}}
\newcommand{\A}{\mathbb{A}}
\newcommand{\PP}{\mathbb{P}}
\newcommand{\E}{\mathbb{E}}
\newcommand{\eps}{\varepsilon}
\newcommand{\dd}{\,\mathrm d}
\newcommand{\sref}[2]{\hyperlink{source:#1:#2}{[S#2]}}
\newif\ifshowcatalogue
\showcataloguetrue
\newcommand{\cataloguescope}[1]{\ifshowcatalogue
 \paragraph{Statement recorded in the supplied source.}
 #1\fi}
\begin{document}
% UnsolvedMath ID 1441; source catalogue code GRAPH-054

\setcounter{section}{1440}

\section{The Second Out-Neighbourhood Conjecture}\label{1441}

{\small\sffamily UnsolvedMath ID 1441 · Graph Theory}

\subsection{Definitions and mathematical statement}

\paragraph{Shared notation.}Unless a section says otherwise, $\N=\{1,2,\ldots\}$,
$\Z$ is the ring of integers, $\Q,\R,\C$ are the rational, real, and complex
fields, $[N]=\{1,\ldots,\lfloor N\rfloor\}$, and $\log$ is the natural
logarithm. For real functions of a parameter tending to infinity,
$f\ll g$ and $f=O(g)$ mean $|f|\leq Cg$ eventually for some fixed $C$;
$f\gg g$ reverses this comparison; $f=o(g)$ means $f/g\to0$;
$f\sim g$ means $f/g\to1$; and $f\asymp g$ means both order bounds.
A subscript on $O$ or $\ll$ specifies allowed constant dependence.
The expression $n^{o(1)}$ in an upper bound means growth smaller than
$n^\epsilon$ for each fixed $\epsilon>0$. Local definitions govern any
exception, finite-field convention, or use of the same symbol for another
object. An asymptotic claim is not automatically an effective algorithm.



An oriented graph has no loops and never has both $u\to v$ and $v\to u$. Let $N^+(v)$ consist of vertices reached by one arc and $N^{++}(v)$ of vertices whose shortest directed distance from $v$ is exactly two. The conjecture asks for a vertex with $|N^{++}(v)|\geq|N^+(v)|$ in every finite oriented graph. \sref{1441}{1} \sref{1441}{2}

\paragraph{Statement recorded in the supplied source.}
Does every oriented graph have a vertex with at least as many vertices at distance 2 as at distance 1? \sref{1441}{1}

\subsection{Short English statement}

Must some vertex in every finite oriented graph reach at least as many vertices for the first time in two steps as it reaches in one step?

\subsection{Research attempt}
\paragraph{Status and convention.}
This is an unresolved attempt at Seymour's conjecture. In particular,
$N^{++}(v)$ excludes both $v$ and $N^+(v)$, even when some out-neighbor
also has a two-step walk from $v$. The calculation below proves the
target when the minimum out-degree is at most two and for orientations
of triangle-free graphs, and gives a general half-size bound. These
are elementary subcases, not improvements on the literature.
The June 2026 preprint of Sadhukhan, Sandeep and Sen reports the case
of minimum out-degree seven, with part of its proof verified by CP-SAT;
its abstract also records the previous threshold six. Those external
results are distinguished from the local proofs and enumerations here.
\url{https://arxiv.org/abs/2606.30588}

\paragraph{A local counting inequality.}
Choose a minimum-out-degree vertex $v$, set $d=d^+(v)$,
$A=N^+(v)$ and $B=N^{++}(v)$, and write $b=|B|$. If $d=0$, the target
holds at $v$. Suppose $d>0$. Every out-neighbor of a vertex $u\in A$
lies in $A\cup B$. It cannot be $v$, because the digraph has no pair
of opposite arcs; if it is outside $A\cup\{v\}$ it lies at distance
exactly two from $v$. Summing the out-degrees of vertices of $A$ gives
\[
 d^2\leq\sum_{u\in A}d^+(u)=e(A)+e(A,B)
       \leq\binom d2+db.
\]
Here $e(A)$ counts arcs internal to $A$ once, and the last inequality
uses at most one arc per unordered pair inside $A$, and at most one
arc from each vertex of $A$ to each vertex of $B$. Thus
\[
 |N^{++}(v)|\geq\left\lceil\frac{d+1}{2}\right\rceil.
\]
For $d=1,2$ this lower bound is $d$, proving the conjecture in those
cases. More precisely, retaining the actual internal arc count yields
$b\geq\lceil d-e(A)/d\rceil$. That refinement identifies the loss in
the argument: out-degree can be absorbed by arcs inside the first
neighborhood instead of producing new second neighbors.

\paragraph{Triangle-free underlying graphs.}
If the underlying undirected graph is triangle-free, there are no arcs
between vertices of $A$: any such edge, together with their edges to
$v$, would form a triangle. Fix $u\in A$. All $d^+(u)$ out-neighbors
then lie in $B$, so $b\geq d^+(u)\geq d$. This proves the target directly,
without regularity or a bound on the number of vertices. It also shows
why bounding internal structure in $A$ can improve the half-size bound.
For an arbitrary oriented graph, however, $A$ may itself be a tournament;
no triangle-free conclusion is available.

\paragraph{Reduction to a terminal strongly connected component.}
Contract the strongly connected components to obtain a finite acyclic
digraph. It has a terminal component $C$ with no arc leaving it. For
every $v\in C$, its first and second out-neighborhoods computed in the
whole digraph are exactly those computed in the induced graph on $C$.
Thus a good vertex of $C$ is good in the whole graph. A singleton
terminal component is a sink and already settles the question.
Consequently it would suffice to prove the conjecture for strongly
connected oriented graphs. This is a valid reduction, but it does not
bound their orders or degrees and is not itself a proof of the target.

\paragraph{Why a degree threshold alone cannot finish the problem.}
Given an oriented graph $G$ and an integer $t\geq1$, replace every vertex
by an independent set of $t$ clones and replace each arc $u\to v$ by
all $t^2$ arcs from the $u$-clones to the $v$-clones. No loops or opposite
arcs appear. For a clone $x$ of $v$, its first neighborhood is the union
of the clone classes of $N^+(v)$, and its second neighborhood is exactly
the union of those of $N^{++}(v)$. In particular,
\[
 d^+(x)=t d^+(v),\qquad |N^{++}(x)|=t|N^{++}(v)|.
\]
There is no new two-step route into $v$'s own clone class: such a route
would require opposite arcs in $G$. A hypothetical counterexample
would therefore produce counterexamples with arbitrarily large minimum
out-degree by increasing $t$. This conditional construction is not a
counterexample; it explains why a fixed verified degree threshold does
not remove the general obstruction.

\paragraph{Exact finite check.}
The companion script enumerates all labeled oriented graphs on one
through five vertices. For each unordered pair it chooses no arc or one
of the two directions, giving $3^{\binom n2}$ graphs at order $n$.
It forms $N^{++}(v)$ as the union of the out-neighborhoods of $N^+(v)$,
then removes $v$ and $N^+(v)$. It checks both the existence of a good
vertex and the displayed bound at every minimum-out-degree vertex.
This is a complete finite enumeration at those orders, including
disconnected graphs, rather than a random search. It offers no
uniform inference for larger orders.

\paragraph{Remaining obstruction.}
The local estimate falls short when many arcs lie inside $A$ or when
arcs from different vertices of $A$ repeatedly hit the same vertices
of $B$. Summing out-degrees counts those repetitions, whereas the
conjecture counts distinct endpoints. A full proof would need a rule
for selecting a different vertex or exploiting these overlaps that
forces $b\geq d$ somewhere. No such rule is established here, and
the conjecture remains unresolved by this attempt.

\subsection{Sources}

{\small

\begin{itemize}

\item[\hypertarget{source:1441:1}{S1}] ULAM.AI, \emph{UnsolvedMath}, supplied \texttt{problems.json}, record 1441 (GRAPH-054), statement and background. The embedded literature-review date is 2026-08-17; that is the archive’s review date, not a fresh certification. Dataset landing page: \url{https://huggingface.co/datasets/ulamai/UnsolvedMath}.

\item[\hypertarget{source:1441:2}{S2}] D. C. Fisher, Squaring a tournament: a proof of Dean's conjecture, Journal of Graph Theory 23 (1996), 43–48. \url{https://doi.org/10.1002/(SICI)1097-0118(199609)23:1%3C43::AID-JGT4%3E3.0.CO;2-K}. Bibliographic information imported from the supplied record.

\item[\hypertarget{source:1441:3}{S3}] Arpan Sadhukhan, R. B. Sandeep, and Sagnik Sen, A proof of Seymour's second neighborhood conjecture for oriented graphs with minimum out-degree equal to 7, arXiv:2606.30588 (2026). \url{https://arxiv.org/abs/2606.30588}. Bibliographic information imported from the supplied record.

\item[\hypertarget{source:1441:4}{S4}] Yandong Bai, Binlong Li, and Boram Park, Towards a strengthening of the second neighborhood conjecture, arXiv:2607.18047 (2026). \url{https://arxiv.org/abs/2607.18047}. Bibliographic information imported from the supplied record.

\end{itemize}

}

\label{end:1441}
\end{document}
